% Created 2026-05-19 Tue 09:50
% Intended LaTeX compiler: xelatex
\documentclass[11pt]{article}
\usepackage{graphicx}
\usepackage{longtable}
\usepackage{wrapfig}
\usepackage{rotating}
\usepackage[normalem]{ulem}
\usepackage{capt-of}
\usepackage{hyperref}
\usepackage{amsmath}
\usepackage{amssymb}
\usepackage{graphicx}
\usepackage{geometry}
\usepackage[a4paper,margin=0.75in, top=0.85in]{geometry}
\usepackage{amsmath,amssymb,mathtools}
\usepackage{microtype}
\usepackage{parskip}
\usepackage{mathpazo}
\usepackage{titlesec}
\titleformat{\section}{\large\bfseries}{}{0em}{}[\titlerule]
\titleformat{\subsection}{\normalize\bfseries}{}{0em}{}
\usepackage{enumitem}
\setlist{noitemsep, topset=4pt}
\usepackage{titling}
\pretitle{\begin{center}\large\bfseries}
\posttitle{\end{center}}
\date{\today}
\title{Homogeneity}
\hypersetup{
 pdfauthor={},
 pdftitle={Homogeneity},
 pdfkeywords={},
 pdfsubject={},
 pdfcreator={Emacs 30.2 (Org mode 9.7.11)}, 
 pdflang={English}}
\begin{document}

\maketitle
\tableofcontents

\section{In First-Order ODEs (Function Scaling)}
\label{sec:org236f0f9}

A firt-order ODE in the form: \(\frac{dy}{dx} = f(x,y)\) is considered homogeneous if the function f(x,y) depends on the ration \(\frac{y}{x}\) .

\textbf{To test for homogeneity:}

$$f(tx, ty) = t^{n} f(x,y)$$

A function is homogeneous of degree \emph{n} if both variables are able to be scaled by constant \emph{t} and yields the original function multiplied by t\textsuperscript{n}.

If the function is rewritten in:

$$M(x,y)dx + N(x,y)dy = 0$$

The function is homogeneous if the \textbf{degree} of \emph{M(x,y)} and \emph{N(x,y)} are equal.

\uline{For example}:

$$(x^{2} + y^{2})dx + xydy = 0$$

This function is homogeneous.

Individually, x\textsuperscript{2} and y\textsuperscript{2} both have degree of 2.

\emph{xy} also has degree of 2. As x\textsuperscript{1} \texttimes{} y\textsuperscript{1} = \emph{xy}.

\textbf{EVERY TERM IS DEGREE 2.}

\begin{itemize}
\item If the degrees are mixed, that is not homogeneous.

\item \textbf{You must also inspect the variables.}

\item If the variables \emph{x} and \emph{y} always appear to be grouped together as a fraction, the function is homogeneous.
\end{itemize}


\uline{Examples}:

\begin{itemize}
\item \(\frac{x}{y}\)

\item \(\frac{y}{x}\)

\item \(\sin\left(\frac{x}{y}\right)\)
\end{itemize}


Homogeneous first-order equations \textbf{cannot have standalone numbers added or subtracted}.
This is becuase constants break the scaling property.

\(\frac{dy}{dx} = \frac{x + y}{x}\) is homogeneous.

\(\frac{dy}{dx} = \frac{x + y + 1}{x}\) is not homogeneous.
\subsection{Solving First-Order ODEs}
\label{sec:org23de1bf}

The idea to solve any first-order ODE is to simplify it by substituting the ratio of x and y.

\(v = \frac{y}{x}\) \(\to\) \(y = vx\)

Remember that:


\begin{align}
v &= \frac{y}{x} \\
y &= vx \\
\frac{dy}{dx} &= v + x\frac{dv}{dx}
\end{align}

Substitute these values into the original equation.

Then to apply \textbf{separation of variables.} Before converting \emph{v} back to the ratio.
\section{Homogeneity in Linear ODEs (Structure and Form)}
\label{sec:org9fa4af8}

A linear ODE is homogeneous if every term contains the dependent variable or one of its derivatives \textbf{AND} if the equation is set to 0.

It can generally be written as:

\[
a_n(x)y^{n} + a_{n-1}(x)y^{n-1}+...+a_1_{}(x)y'+ a_0_{}(x)y = 0
\]

There can be \textbf{no forcing term}.

If \emph{y = 0} is a valid solution, it is homogeneous.

If y\textsubscript{1} and y\textsubscript{2} are solutions to a linear homogeneous ODE, then any linear combination of \(C_1 y_1 + C_2 y_2\) is also a solution.

\textbf{To solve}

If the coefficients are constant, use a characteristic polynomial.
\end{document}
